\begin{frame}[fragile]{Curto-circuito: nem sempre avaliar tudo}
\begin{lstlisting}
int quantidade = 0;
int total = 100;
bool media_alta = (quantidade != 0)
               && (total / quantidade > 20);
\end{lstlisting}
\begin{itemize}
\item \texttt{\&\&} avalia primeiro a esquerda. Se ela for falsa, não avalia a direita.
\item No exemplo, a divisão não ocorre; \texttt{media\_alta} é falso.
\item \texttt{||} também começa à esquerda: se ela for verdadeira, não avalia a direita.
\end{itemize}
\begin{block}{A ordem protege a operação}
Verificar o divisor deve vir antes da divisão. Inverter as duas expressões elimina essa proteção.
\end{block}
\end{frame}
